\documentclass[11pt,a4paper]{article}

% --- Paquetes ---
\usepackage[utf8]{inputenc}
\usepackage[T1]{fontenc}
\usepackage[spanish]{babel}
\usepackage{amsmath, amssymb, amsthm}
\usepackage{geometry}
\usepackage{graphicx}
\usepackage{hyperref}
\usepackage{xcolor}
\usepackage{listings}
\usepackage{booktabs}
\usepackage{enumitem}
\usepackage{fancyhdr}
\usepackage{titlesec}

% --- Configuracion de pagina ---
\geometry{left=2.5cm, right=2.5cm, top=2.5cm, bottom=2.5cm}
\setlength{\parindent}{0pt}
\setlength{\parskip}{0.6em}

% --- Colores ---
\definecolor{azulprofundo}{RGB}{30, 58, 138}
\definecolor{azulclaro}{RGB}{59, 130, 246}
\definecolor{grisclaro}{RGB}{240, 240, 245}

% --- Hyperref ---
\hypersetup{
    colorlinks=true,
    linkcolor=azulprofundo,
    urlcolor=azulclaro,
    citecolor=azulprofundo
}

% --- Titulos de seccion ---
\titleformat{\section}{\Large\bfseries\color{azulprofundo}}{\thesection}{1em}{}
\titleformat{\subsection}{\large\bfseries\color{azulclaro}}{\thesubsection}{1em}{}

% --- Header/Footer ---
\pagestyle{fancy}
\fancyhf{}
\rhead{\small Sistemas Lineales y Senales}
\lhead{\small Restauracion de Placas Vehiculares}
\cfoot{\thepage}

% --- Codigo ---
\lstset{
    basicstyle=\ttfamily\small,
    backgroundcolor=\color{grisclaro},
    frame=single,
    framesep=6pt,
    rulecolor=\color{gray!30},
    breaklines=true,
    numbers=left,
    numberstyle=\tiny\color{gray},
    keywordstyle=\color{azulprofundo}\bfseries,
    commentstyle=\color{green!50!black},
    stringstyle=\color{red!70!black},
    showstringspaces=false,
    tabsize=4
}

% =========================================================
\begin{document}

% --- Portada ---
\begin{titlepage}
    \centering
    \vspace*{2cm}

    {\LARGE\bfseries\color{azulprofundo} Restauracion de Imagenes \\[0.3em]
    de Placas Vehiculares}\\[0.8em]
    {\Large mediante Convolucion y Deconvolucion}\\[2.5cm]

    {\large\textit{Propuesta de Proyecto Final}}\\[0.5em]
    {\large\textbf{Sistemas Lineales y Senales}}\\[3cm]

    \rule{0.6\textwidth}{0.5pt}\\[0.8em]
    {\large Aplicacion de sistemas LTI al procesamiento de imagenes}\\
    \rule{0.6\textwidth}{0.5pt}\\[3cm]

    \begin{tabular}{rl}
        \textbf{Realizado por:} & \rule{6cm}{0.4pt} \\
        \textbf{Docente:}       & \rule{6cm}{0.4pt} \\
        \textbf{Fecha:}         & \today \\
    \end{tabular}

    \vfill
\end{titlepage}

% =========================================================
\section{Resumen}

Se propone desarrollar un sistema computacional que demuestre, de forma matematica y practica, como la degradacion de imagenes capturadas por camaras de transito puede modelarse como un \textbf{sistema lineal e invariante en el tiempo} (LTI) mediante la operacion de \textbf{convolucion 2D}, y como dicha degradacion puede revertirse parcialmente mediante tecnicas de \textbf{deconvolucion}.

El proyecto tendra como caso de estudio la restauracion de \textbf{placas vehiculares}, con aplicacion directa en control de transito, peajes automaticos, radares y sistemas de videovigilancia.

% =========================================================
\section{Planteamiento del Problema}

Las camaras de transito, radares y sistemas de vigilancia frecuentemente capturan imagenes de placas vehiculares en condiciones adversas: movimiento del vehiculo, desenfoque optico, condiciones de baja luz o ruido del sensor. Estas degradaciones comprometen la legibilidad de la placa y afectan directamente la precision de los sistemas automaticos de reconocimiento (ALPR, \textit{Automatic License Plate Recognition}).

Se plantea abordar este problema desde los principios fundamentales de los \textbf{sistemas lineales}: modelar la camara como un sistema LTI, caracterizarlo por su respuesta al impulso, y aplicar tecnicas matematicas de inversion para recuperar la imagen original.

% =========================================================
\section{Objetivos}

\subsection{Objetivo General}
Desarrollar un sistema de procesamiento de imagenes que modele la degradacion como convolucion e implemente metodos de deconvolucion para restaurar la legibilidad de placas vehiculares, validando los resultados mediante OCR.

\subsection{Objetivos Especificos}
\begin{enumerate}[leftmargin=*]
    \item Implementar manualmente el operador de convolucion 2D con bucles explicitos, demostrando la aplicacion directa de la definicion matematica.
    \item Construir kernels que modelen degradaciones realistas (desenfoque Gaussiano, desenfoque por movimiento).
    \item Simular el proceso de degradacion combinando convolucion y ruido aditivo.
    \item Implementar tres estrategias de deconvolucion: filtro inverso, filtro de Wiener y Richardson-Lucy iterativo.
    \item Evaluar cuantitativamente la calidad de la restauracion mediante metricas objetivas (MSE, PSNR, SSIM).
    \item Validar de forma practica la mejora mediante un motor OCR que lea las placas antes y despues de la restauracion.
\end{enumerate}

% =========================================================
\section{Marco Teorico}

\subsection{Sistema Lineal e Invariante en el Tiempo}

Un sistema $T\{\cdot\}$ que toma una senal $f(x,y)$ y produce una senal $g(x,y)$ es \textbf{LTI} si cumple:

\begin{align}
\text{\textbf{Linealidad:}} \quad &T\{a\,f_1 + b\,f_2\} = a\,T\{f_1\} + b\,T\{f_2\} \\
\text{\textbf{Invariancia:}} \quad &T\{f(x-x_0, y-y_0)\} = g(x-x_0, y-y_0)
\end{align}

Un sistema LTI queda completamente caracterizado por su \textbf{respuesta al impulso} $h(x,y)$, tambien conocida como PSF (\textit{Point Spread Function}):

\begin{equation}
h(x,y) = T\{\delta(x,y)\}
\end{equation}

donde $\delta(x,y)$ es la funcion delta de Dirac bidimensional.

\subsection{Modelo Matematico de Degradacion}

La respuesta de un sistema LTI a cualquier senal de entrada es la \textbf{convolucion} de la entrada con la respuesta al impulso. En presencia de ruido aditivo, el modelo de degradacion de la imagen es:

\begin{equation}
\boxed{\; g(x,y) = f(x,y) * h(x,y) + n(x,y) \;}
\end{equation}

donde:
\begin{itemize}[leftmargin=*, noitemsep]
    \item $f(x,y)$: imagen original (senal de entrada limpia)
    \item $h(x,y)$: PSF del sistema (kernel de degradacion)
    \item $n(x,y)$: ruido aditivo (tipicamente Gaussiano, $n \sim \mathcal{N}(0, \sigma^2)$)
    \item $g(x,y)$: imagen degradada observada
    \item $*$: operador de convolucion 2D
\end{itemize}

\subsection{Convolucion 2D Discreta}

La convolucion 2D discreta entre una imagen $f[m,n]$ y un kernel $h[i,j]$ se define como:

\begin{equation}
\boxed{\; g[m,n] = \sum_{i} \sum_{j} f[m-i,\, n-j] \cdot h[i,j] \;}
\end{equation}

Esta es la piedra angular del proyecto y se implementara \textbf{manualmente con bucles explicitos}, evidenciando que cada pixel de salida es una suma ponderada de los vecinos del pixel de entrada, donde los pesos estan dados por el kernel volteado.

\subsection{Dominio de la Frecuencia}

Por el \textbf{Teorema de Convolucion}, la convolucion en el dominio espacial equivale a la multiplicacion punto a punto en el dominio de la frecuencia:

\begin{equation}
\mathcal{F}\{f * h\} = F(u,v) \cdot H(u,v)
\end{equation}

El modelo de degradacion en frecuencia se escribe como:

\begin{equation}
G(u,v) = F(u,v) \cdot H(u,v) + N(u,v)
\end{equation}

Esta representacion permite atacar el problema de deconvolucion mediante operaciones algebraicas sobre matrices complejas de frecuencia.

\subsection{Estrategias de Deconvolucion}

\subsubsection{Filtro Inverso}

La solucion mas directa consiste en dividir el espectro observado por la PSF:

\begin{equation}
\hat{F}(u,v) = \frac{G(u,v)}{H(u,v)}
\end{equation}

\textbf{Limitacion:} en frecuencias donde $|H(u,v)| \to 0$, el ruido $N/H$ se amplifica catastroficamente. Se incluira con proposito \textbf{didactico} para evidenciar la necesidad de regularizacion.

\subsubsection{Filtro de Wiener}

Introduce un parametro de regularizacion $K$ que aproxima la relacion inversa senal-ruido:

\begin{equation}
\boxed{\; \hat{F}(u,v) = \left[ \frac{H^*(u,v)}{|H(u,v)|^2 + K} \right] \cdot G(u,v) \;}
\end{equation}

donde $H^*$ es el conjugado complejo de $H$. Sera el \textbf{metodo principal} del proyecto por su robustez y fundamento estadistico (minimiza el error cuadratico medio esperado).

\subsubsection{Richardson-Lucy Iterativo}

Metodo iterativo derivado de la estadistica bayesiana que preserva la positividad:

\begin{equation}
f_{k+1}(x,y) = f_k(x,y) \cdot \left[ \frac{g(x,y)}{f_k(x,y) * h(x,y)} * h^T(x,y) \right]
\end{equation}

donde $h^T$ es la PSF transpuesta (rotada 180 grados). Converge a la estimacion de maxima verosimilitud.

\subsection{Metricas de Calidad Objetiva}

\begin{align}
\text{MSE} &= \frac{1}{MN} \sum_{x,y} [f(x,y) - \hat{f}(x,y)]^2 \\[0.3em]
\text{PSNR} &= 10 \cdot \log_{10}\left( \frac{\text{MAX}^2}{\text{MSE}} \right) \quad \text{[dB]} \\[0.3em]
\text{SSIM}(x,y) &= \frac{(2\mu_x\mu_y + C_1)(2\sigma_{xy} + C_2)}{(\mu_x^2 + \mu_y^2 + C_1)(\sigma_x^2 + \sigma_y^2 + C_2)}
\end{align}

% =========================================================
\section{Metodologia}

El proyecto se desarrollara en seis fases secuenciales:

\begin{enumerate}[leftmargin=*]
    \item \textbf{Implementacion del nucleo matematico:} convolucion 2D manual con bucles explicitos, manejo de padding y diferentes modos de contorno.
    \item \textbf{Construccion de kernels:} Gaussiano, motion blur, Laplaciano y Sobel, todos construidos a partir de sus formulas matematicas explicitas.
    \item \textbf{Simulacion de degradacion:} aplicacion del modelo $g = f * h + n$ con diferentes configuraciones de ruido.
    \item \textbf{Implementacion de deconvolucion:} los tres metodos descritos, operando sobre la representacion en frecuencia obtenida via DFT.
    \item \textbf{Evaluacion cuantitativa:} calculo de metricas MSE, PSNR y SSIM comparando cada etapa.
    \item \textbf{Validacion practica:} ejecucion de OCR sobre la imagen original, degradada y restaurada para cuantificar la mejora en legibilidad.
\end{enumerate}

% =========================================================
\section{Tecnologias}

\begin{table}[h]
\centering
\begin{tabular}{@{}lll@{}}
\toprule
\textbf{Componente} & \textbf{Tecnologia} & \textbf{Proposito} \\
\midrule
Lenguaje & Python 3.8+ & Desarrollo principal \\
Computo numerico & NumPy & Operaciones matriciales, FFT \\
Procesamiento imagen & OpenCV & Carga de imagenes, CLAHE \\
Visualizacion & Matplotlib & Graficas y espectros \\
Interfaz grafica & CustomTkinter & GUI moderna con tema oscuro \\
OCR & RapidOCR (ONNX Runtime) & Lectura de placas (PaddleOCR) \\
Validacion & scikit-image & Verificacion de metricas \\
\bottomrule
\end{tabular}
\caption{Tecnologias propuestas para la implementacion.}
\end{table}

\textbf{Restriccion fundamental:} las operaciones matematicas centrales (convolucion 2D y logica de deconvolucion) seran implementadas manualmente sin delegar en funciones de libreria. Las herramientas externas se usaran unicamente para:
\begin{itemize}[leftmargin=*, noitemsep]
    \item Entrada/salida de imagenes (OpenCV)
    \item Calculo numerico de la DFT (NumPy FFT, un algoritmo, no una funcion de restauracion)
    \item Visualizacion e interfaz
    \item Validacion de metricas propias
\end{itemize}

% =========================================================
\section{Arquitectura Propuesta}

La solucion se organizara en modulos independientes bajo una carpeta \texttt{src/}:

\begin{lstlisting}[language=,numbers=none]
proyecto_sistemas/
|-- app.py              # Interfaz grafica (CustomTkinter)
|-- main.py             # Pipeline de consola
|-- requirements.txt
|-- images/             # Imagenes de prueba
|-- output/             # Resultados generados
`-- src/
    |-- convolution.py   # NUCLEO: convolucion 2D manual
    |-- kernels.py       # Kernels matematicos
    |-- degradation.py   # Modelo g = f*h + n
    |-- restoration.py   # Deconvolucion: Wiener, R-L, inverso
    |-- postprocessing.py # Sharpening, CLAHE, bordes
    |-- ocr.py           # Motor de reconocimiento
    |-- metrics.py       # MSE, PSNR, SSIM manuales
    `-- visualization.py
\end{lstlisting}

\subsection{Fragmento clave: convolucion manual}

El corazon matematico del proyecto sera un doble bucle que implementa directamente la definicion de convolucion:

\begin{lstlisting}[language=Python]
def convolve2d_manual(image, kernel, padding='zero'):
    # Voltear el kernel (definicion estricta de convolucion)
    kernel_flipped = kernel[::-1, ::-1]

    kH, kW = kernel_flipped.shape
    padded = pad_image(image, kH // 2, kW // 2, mode=padding)

    M, N = image.shape
    output = np.zeros((M, N))

    for m in range(M):
        for n in range(N):
            region = padded[m:m + kH, n:n + kW]
            output[m, n] = np.sum(region * kernel_flipped)

    return output
\end{lstlisting}

Cada pixel de salida es una suma ponderada de los vecinos del pixel de entrada, demostrando la aplicacion directa de la formula matematica.

% =========================================================
\section{Entregables Esperados}

\begin{itemize}[leftmargin=*]
    \item Codigo fuente modular, documentado con las formulas matematicas correspondientes.
    \item Interfaz grafica que permita cargar imagenes, ajustar parametros de degradacion y restauracion, y visualizar resultados comparativos.
    \item Visualizaciones: espectros de Fourier, mapas de diferencia, graficas comparativas, tabla de metricas y resultados OCR anotados.
    \item Reporte de resultados cuantitativos para al menos una imagen de prueba, mostrando la mejora en PSNR, SSIM y precision del OCR.
    \item Documentacion tecnica (README) con instrucciones de instalacion y uso.
\end{itemize}

% =========================================================
\section{Aplicaciones Practicas}

El sistema resultante sera aplicable en los siguientes dominios:

\begin{itemize}[leftmargin=*, noitemsep]
    \item Control de transito y radares automaticos de velocidad
    \item Peajes electronicos con reconocimiento de placas
    \item Sistemas de videovigilancia y seguridad vial
    \item Lectura automatica de placas (ALPR) en estacionamientos
    \item Investigacion forense en casos de imagenes deterioradas
\end{itemize}

% =========================================================
\section{Conclusion}

Este proyecto vinculara los conceptos teoricos del curso de \textit{Sistemas Lineales y Senales} con una aplicacion tangible y de valor practico. La implementacion manual de la convolucion y la deconvolucion evidenciara que los fundamentos matematicos no son abstracciones alejadas de la realidad, sino herramientas concretas que permiten resolver problemas reales en el procesamiento digital de imagenes.

Se demostrara que un sistema fisico (una camara) puede modelarse como un sistema LTI, caracterizarse matematicamente, y su efecto revertirse mediante operaciones algebraicas fundamentadas en la teoria de senales.

% =========================================================
\section*{Referencias}
\addcontentsline{toc}{section}{Referencias}

\begin{itemize}[leftmargin=*, noitemsep]
    \item Gonzalez, R.C. \& Woods, R.E. (2018). \textit{Digital Image Processing} (4th ed.). Pearson.
    \item Oppenheim, A.V. \& Willsky, A.S. (1997). \textit{Signals and Systems} (2nd ed.). Prentice Hall.
    \item Wang, Z., Bovik, A.C., Sheikh, H.R. \& Simoncelli, E.P. (2004). Image Quality Assessment: From Error Visibility to Structural Similarity. \textit{IEEE Transactions on Image Processing}, 13(4), 600--612.
    \item Richardson, W.H. (1972). Bayesian-Based Iterative Method of Image Restoration. \textit{Journal of the Optical Society of America}, 62(1), 55--59.
    \item Lucy, L.B. (1974). An Iterative Technique for the Rectification of Observed Distributions. \textit{The Astronomical Journal}, 79, 745--754.
\end{itemize}

\end{document}
